\subsection{Compact linear maps and weak topologies}

Let E be a locally convex space, \(E'\) its dual. Recall that the topology \(\sigma(\mathrm{E},\mathrm{E}')\) on E is called the weak topology on E (TVS, IV, p. 4). In this section, \(E_{\sigma}\) denotes the space E equipped with the weak topology.

PROPOSITION 6. --- Let E be a locally convex space, F a separated locally convex space, \(u: \mathrm{E} \to \mathrm{F}\) a compact linear map and B a bounded subset of E. The restriction of \(u\) to B is a continuous map from the set B, equipped with the topology induced by \(\sigma(\mathrm{E}, \mathrm{E}')\) , into the space F.

The map \(u\) is a continuous linear map from E into F, and also from \(\mathrm{E}_{\sigma}\) into \(\mathrm{F}_{\sigma}\) . Its restriction to B is therefore a continuous map from the set B, equipped with the topology induced by \(\sigma(\mathrm{E}, \mathrm{E}')\) , into the space \(\mathrm{F}_{\sigma}\) . Now the set \(u(\mathrm{B})\) is contained in a compact subset C of F (remark 1 of III, p. 2), and the space \(\mathrm{F}_{\sigma}\) is Hausdorff, so that the induced topologies on C from those of F and of \(\mathrm{F}_{\sigma}\) coincide. The proposition follows.

COROLLARY. --- Let \((x_{n})_{n\in\mathbf{N}}\) be a sequence of points of E, which converges to a point x of E for the weak topology. The sequence \((u(x_{n}))_{n\in\mathbf{N}}\) converges in F to \(u(x)\) .

Indeed, the set consisting of the point \(x\) and the points of the sequence \((x_{n})_{n\in \mathbf{N}}\) is a bounded subset of E (EVT, III, p. 3, cor. of prop. 2).

PROPOSITION 7. --- Let E be a semi-normed space, E' its dual, B the unit ball of E', and u a linear map from E' into a separated locally convex space F. If the restriction of u to B equipped with the topology induced by \(\sigma(\mathrm{E}',\mathrm{E})\) is continuous, then \(u(\mathrm{B})\) is a compact subset of F and u is a compact linear map from the strong dual \(\mathrm{E}_b'\) of E into F.

Indeed, the set B, equipped with the topology induced by \(\sigma(\mathrm{E}^{\prime},\mathrm{E})\) is compact (EVT, III, p. 17, cor. 2).

COROLLARY 1. --- Let E be a semi-normed space of countable type, let E' be its dual and let B be the unit ball of E'. Let u be a linear mapping from E' into a separated locally convex space F. Suppose that for every sequence \((x_{n})_{n\in\mathbf{N}}\) of elements of B that converges for \(\sigma(\mathrm{E}',\mathrm{E})\) to 0, the sequence \((u(x_{n}))_{n\in\mathbf{N}}\) converges to 0 in F. Then, \(u(\mathrm{B})\) is a compact subset of F and u is a compact linear map from the strong dual \(E_{b}'\) of E into F.

Indeed, the topology induced on B by \(\sigma(\mathrm{E}^{\prime},\mathrm{E})\) is metrizable (EVT, III, p. 19, cor. 2). The hypothesis of the corollary therefore implies that the restriction of \(u\) to B is a continuous map from B into F, when B is equipped with the topology \(\sigma(\mathrm{E}^{\prime},\mathrm{E})\) and the corollary follows from the preceding proposition.

COROLLARY 2. --- Let E be a semi-normed space and S a set of precompact subsets of E whose union is dense in E. Denote by \(E'_{b}\) the strong dual of E and by \(E'_{S}\) the dual of E equipped with the S-topology. The space \(E'_{S}\) is separated and the identity map of \(E'_{b}\) into \(E'_{S}\) is compact.

The space \(E_{S}^{\prime}\) is separated according to TG, X, p. 8, prop. 7; on the unit ball of \(E^{\prime}\) the S-topology coincides with the weak topology \(\sigma(\mathrm{E}^{\prime},\mathrm{E})\) (EVT, III, p. 17, prop. 5), whence the corollary.

Example. --- Let E be a semi-normed space. The identity map of \(E_{b}^{\prime}\) into the space \(E^{\prime}\) , endowed with the weak topology, the topology of compact convergence, or the topology of precompact convergence, is compact.

PROPOSITION 8. --- Let E be a reflexive Banach space, B its unit ball, F a separated locally convex space, and u: E → F a linear map. The following conditions are equivalent:

\begin{enumerate}
\def\labelenumi{(\roman{enumi})}
\item
  The linear map u is compact;
\item
  The set \(u(\mathrm{B})\) is a compact subset of F;
\item
  The restriction of u to B is a continuous map from the set B, equipped with the topology induced by \(\sigma(\mathrm{E},\mathrm{E}')\) into F;
\item
  For any sequence \((x_{n})_{n\in \mathbf{N}}\) of points of B that converges to 0 for the topology \(\sigma (\mathrm{E},\mathrm{E}^{\prime})\) the sequence \((u(x_n))_{n\in \mathbf{N}}\) converges to 0 in F;
\item
  From any infinite sequence of points of \(u(\mathrm{B})\) one can extract a subsequence that converges in F;
\item
  Any infinite sequence of points of \(u(\mathrm{B})\) has a cluster point in F.
\end{enumerate}

Since E is a reflexive space, B is a compact subset of \(E_{\sigma}\) (EVT, IV, p. 17, prop. 6), hence (iii) implies (ii). Condition (ii) implies (i) by definition, and condition (i) implies (iii) by prop. 6. It is also elementary that (iii) implies (iv) and that (v) implies (vi).

We will now show that (iv) implies (v). Let \((x_{n})\) be an infinite sequence of points of B. Since B is a compact subset of \(E_{\sigma}\) , by Šmulian's theorem (EVT, IV, p. 36, th. 2) there exists a sequence \((y_{n})\) , extracted from the sequence \((x_{n})\) , which converges in \(E_{\sigma}\) to a limit y. The sequence \((y_{n}-y)\) is bounded in \(E_{\sigma}\) , hence in E (EVT, IV, p. 1, prop. 1); it is therefore contained in a homothetic set of B. If condition (iv) is satisfied, the sequence \((u(y_{n}-y))\) converges to 0 in F, and \((u(y_{n}))\) , which is a subsequence of \((u(x_{n}))\) , converges to \(u(y)\) . This proves that (iv) implies (v).

Finally, let us show that (vi) implies (ii). Let \(j: \mathrm{F} \to \widehat{\mathrm{F}}\) be the canonical injection of F into its completion. Under hypothesis (vi), \(u(\mathrm{B})\) is a precompact subset of F (EVT, IV, p. 32, prop. 1), hence \(j(u(\mathrm{B}))\) is a relatively compact subset of \(\widehat{\mathrm{F}}\) (TG, II, no. 2, p. 29) and \(j \circ u\) is a compact linear map. By the already proven equivalence of conditions (i) and (ii), \(j(u(\mathrm{B}))\) is a compact subset of \(\widehat{\mathrm{F}}\) . Consequently, \(u(\mathrm{B})\) is a compact subset of F.

